\section{Skills 概述}

Skills 已经成为 Claude Code 中使用最广泛的\textbf{扩展点}（extension point）之一。它们灵活、容易制作，分发起来也很简单。但也正因为太灵活，实践中会面临三个核心问题：
\begin{enumerate}
    \item 什么类型的 Skills 值得做？
    \item 写出好 Skill 的秘诀是什么？
    \item 什么时候该把它们分享给别人？
\end{enumerate}

Anthropic 内部已有数百个 Skills 在活跃使用。本文就是他们从中总结出的经验教训。

\begin{importantbox}{Skills 不只是 Markdown 文件}
一个常见误解是认为 Skills ``只不过是 markdown 文件''。实际上，\textbf{Skills 最有趣的地方恰恰在于它们不只是文本文件}。Skills 是\textbf{文件夹}，可以包含脚本、资源文件（assets）、数据等内容，智能体可以发现、探索和操作这些内容。在 Claude Code 中，Skills 还拥有丰富的配置选项，包括注册动态钩子（hooks）。Anthropic 发现，最有意思的 Skills 往往创造性地利用了这些配置选项和文件夹结构。
\end{importantbox}

\subsection{本章小结}

Skills 是 Claude Code 的核心扩展机制，其本质是包含指令、脚本和资源的文件夹，而非简单的文本文件。理解这一点是用好 Skills 的前提。

